% QED Limit Bridge — Cross-Reference Summary
% Version: 1.0
% Date: 2026-03-10
% Status: Canonical bridge — cross-references only; no new derivations
% Author: Ing. David Jaroš
% © 2026 Ing. David Jaroš — CC BY-NC-ND 4.0
%
% PURPOSE: This file collects the QED limit claims and locates them
% across the repository, with explicit status labels (proved /
% semi-empirical / open) so that an external reviewer can assess
% the QED sector without reading multiple scattered files.

\documentclass[11pt,a4paper]{article}
\usepackage{amsmath,amssymb,amsthm}
\usepackage{hyperref}
\usepackage{geometry}\geometry{margin=2.5cm}
\usepackage{xcolor}
\usepackage{tcolorbox}\tcbuselibrary{skins,breakable}
\usepackage{booktabs}
\newcommand{\Lzero}{\textbf{[L0]}}
\newcommand{\Lone}{\textbf{[L1]}}
\newcommand{\Ltwo}{\textbf{[L2]}}
\newcommand{\OPEN}{\textbf{[OPEN]}}
\newcommand{\SE}{\textbf{[SE]}}
\begin{document}

\section{QED Limit: Navigation Guide and Status}
\label{sec:bridges:qed_limit}

\begin{tcolorbox}[colback=green!5!white,colframe=green!50!black,title=QED Limit Status Summary]
\textbf{PROVED:} QED Lagrangian, U(1) gauge invariance, massless photon, $B_0 = 8\pi$ \\
\textbf{PROVED:} Running coupling form $1/\alpha(\mu) = 1/\alpha(\mu_0) + (B_\alpha/2\pi)\ln(\mu/\mu_0)$ \\
\textbf{SEMI-EMPIRICAL:} Precise value $B_\alpha \approx 46.3$ (requires $B_\text{base}$, which is open) \\
\textbf{OPEN HARD PROBLEM:} $B_\text{base} = N_\text{eff}^{3/2} = 41.57$ — exponent 3/2 not derived \\
\textbf{OPEN HARD PROBLEM:} Correction factor $R \approx 1.114$ — geometric origin unknown
\end{tcolorbox}

\begin{tcolorbox}[colback=blue!5!white,colframe=blue!40!black,
  title=Status Update (v44)]
\textbf{B\_base exponent $3/2$}: now \textbf{Partially Proved [L1]} via
T³ heat kernel (see \texttt{canonical/interactions/B\_base\_derivation\_complete.tex}).
Ab initio beta-fn coefficient matching remains \textbf{[Open]}.\\[4pt]
\textbf{R $\approx$ 1.114}: remains \textbf{Open Hard Problem}.
\end{tcolorbox}

\subsection{QED Lagrangian (PROVED)}

\textbf{Claim:} The QED Lagrangian emerges from the UBT kinetic term via U(1) gauge coupling:
\[
\mathcal{L}_\text{QED} = \mathrm{Tr}[(D_\mu\Theta)^\dagger(D^\mu\Theta)] - \tfrac{1}{4}F_{\mu\nu}F^{\mu\nu},
\qquad D_\mu = \partial_\mu + igA_\mu.
\]

\textbf{Location:} \texttt{canonical/interactions/qed.tex},
Section~\ref{sec:canonical:qed}. U(1) gauge invariance
($\Theta \to e^{i\alpha(x)}\Theta$, $A_\mu \to A_\mu - g^{-1}\partial_\mu\alpha$)
and gauge fixing are canonical there.

\subsection{Maxwell Limit (PRESENT; follows from QED Lagrangian)}

\textbf{Claim:} The Maxwell kinetic term $-\tfrac{1}{4}F_{\mu\nu}F^{\mu\nu}$ with
$F_{\mu\nu} = \partial_\mu A_\nu - \partial_\nu A_\mu$ is the Maxwell Lagrangian.
The Maxwell limit is obtained by setting matter fields to their vacuum values
(constant phase $\varphi = \mathrm{const}$).

\textbf{Location:} \texttt{canonical/interactions/qed.tex} — $F_{\mu\nu}$ definition.
No dedicated Maxwell derivation file exists; the limit is immediate from the
QED Lagrangian.

The gauge field contribution to $T_{\mu\nu}$ is derived in
\texttt{docs/papers/appendix\_GR\_embedding.tex}.

\subsection{Massless Photon (PROVED)}

\textbf{Claim:} The photon is massless at constant phase $\varphi = \mathrm{const}$
(no Proca mass term generated).

\textbf{Location:} \texttt{canonical/interactions/qed.tex};
\texttt{DERIVATION\_INDEX.md}: ``massless photon: verified''.

\subsection{\texorpdfstring{$B_0 = 8\pi$ from $S_\text{kin}[\Theta]$}{B0 = 8pi} (PROVED [L1])}

\textbf{Claim:} The one-loop baseline $B_0 = 8\pi = 2\pi N_\text{eff}/3$
is a zero-free-parameter result from the kinetic action $S_\text{kin}[\Theta]$.
The QED limit $N_\text{eff} = 1$ gives $B_0 = 2\pi/3$, consistent.

\textbf{Location:} \texttt{STATUS\_ALPHA.md}, §5;
\texttt{consolidation\_project/N\_eff\_derivation/step2\_vacuum\_polarization.tex}
(Theorem 3.1);
\texttt{DERIVATION\_INDEX.md}: ``$B_0 = 25.1$ (one-loop baseline): Proved [L1]''.

\subsection{Running Coupling \texorpdfstring{$\alpha(\mu)$}{alpha(mu)} (PROVED structure; SEMI-EMPIRICAL numerics)}

\textbf{Claim:} The running coupling formula
\[
\frac{1}{\alpha(\mu)} = \frac{1}{\alpha(\mu_0)} + \frac{B_\alpha}{2\pi}\ln\frac{\mu}{\mu_0}
\]
is derived from $N_\text{eff} = 12$. The \emph{structural} form is proved;
the precise coefficient $B_\alpha \approx 46.3$ depends on the open $B_\text{base}$ value.

\textbf{Location:} \texttt{STATUS\_ALPHA.md}, §7.

\subsection{\texorpdfstring{$B_\text{base} = N_\text{eff}^{3/2} = 41.57$}{B\_base} (OPEN HARD PROBLEM)}

\begin{tcolorbox}[colback=orange!10!white,colframe=orange!75!black,title=Open Hard Problem: $B_\text{base}$]
The formula $B_\text{base} = N_\text{eff}^{3/2} = 12^{3/2} = 41.57$ is a
\textbf{motivated conjecture}, not a first-principles derivation.

\textbf{What is proved:}
\begin{itemize}
  \item Factor 3 = $\dim_\mathbb{R}(\mathrm{Im}\,\mathbb{H})$ — algebraic fact [L0]
  \item Factor 2 from Gaussian integration — proved [L1]
\end{itemize}

\textbf{What remains open:}
\begin{itemize}
  \item The exponent 3/2 (22 approaches tested; all dead ends except
        Kac-Moody $k=1$ motivated conjecture via H2)
  \item Gap G3-k: direct computation of Kac-Moody level $k$ from $S[\Theta]$
  \item Gap G8: modular weight of UBT partition function $\hat{Z}(\tau)$
\end{itemize}

\textbf{Do not invent a derivation.} The honest status is documented in
\texttt{STATUS\_ALPHA.md}, §9 and
\texttt{DERIVATION\_INDEX.md} (row: $B_\text{base} = N_\text{eff}^{3/2}$).
\end{tcolorbox}

\subsection{Correction Factor \texorpdfstring{$R \approx 1.114$}{R} (OPEN HARD PROBLEM)}

The correction factor bridging $B_\text{base} = 41.57 \to B_\alpha \approx 46.3$
has unknown geometric origin. Eight approaches tested; all dead ends.
See \texttt{docs/PROOFKIT\_ALPHA.md}, §5;
\texttt{consolidation\_project/alpha\_derivation/r\_factor\_geometry.tex}.

\subsection{Notation disambiguation: \texorpdfstring{$B_\text{base}$ vs.\ $B_m$}{B\_base vs B\_m}}

\begin{tcolorbox}[colback=blue!5!white,colframe=blue!40!black,title=Notation note]
$B_\text{base} \approx 41.57$ (dimensionless) is the running-coupling coefficient
used in the $\alpha$ derivation.

$B_m \approx -14.099\ \mathrm{MeV}$ is the logarithmic correction in the Hopfion
lepton mass formula $m(n) = A\,n^p - B_m\,n\ln n$.

These are \emph{distinct} objects with different dimensions and physical interpretations.
\end{tcolorbox}

\subsection{Further reading}

\begin{itemize}
  \item \texttt{canonical/interactions/qed.tex} — QED Lagrangian (canonical)
  \item \texttt{STATUS\_ALPHA.md} — full $\alpha$ derivation chain and gap list
  \item \texttt{DERIVATION\_INDEX.md} — row-by-row derivation status for all $\alpha$ steps
  \item \texttt{consolidation\_project/alpha\_derivation/} — all B\_base attempts
  \item \texttt{docs/papers/appendix\_GR\_embedding.tex} — gauge field contribution to $T_{\mu\nu}$
\end{itemize}

% End of QED_limit_bridge.tex

\end{document}
